\section{Complete Final Outcomes by Training Seed}
%% ============================================================

\begin{table}[t]
\centering
\caption{Final evaluation of every completed PPO policy pair. Each row averages 64 evaluation episodes. B denotes informative baseline and A denotes action-independent retraining. Seed suffix 01 denotes 20261001, continuing through 10. Rank columns describe projected, time-augmented encoder features.}
\label{tab:seed_results}
\small
\sisetup{table-format=1.6, round-mode=places, round-precision=6}
\begin{tabular}{@{}ccSSSSS@{}}
\toprule
Seed & Arm & {Price} & {Elevation} & {Surplus} & {Rank 0} & {Rank 1} \\
\midrule
01 & B & 1.497102 & 0.053323 & 1.157867 & 1.276308 & 1.237626 \\
01 & A & 1.488946 & 0.035278 & 1.159020 & 1.418177 & 1.386605 \\
02 & B & 1.481101 & 0.017922 & 1.159558 & 1.218294 & 1.241206 \\
02 & A & 1.476651 & 0.008077 & 1.160178 & 1.547627 & 1.522354 \\
03 & B & 1.468815 & -0.009260 & 1.161126 & 1.229325 & 1.311140 \\
03 & A & 1.458835 & -0.031339 & 1.162020 & 1.328914 & 1.342101 \\
04 & B & 1.479467 & 0.014307 & 1.160148 & 1.254956 & 1.291669 \\
04 & A & 1.477020 & 0.008893 & 1.160382 & 1.232180 & 1.275160 \\
05 & B & 1.509558 & 0.080881 & 1.155365 & 1.281766 & 1.237738 \\
05 & A & 1.466939 & -0.013408 & 1.161483 & 1.331038 & 1.422554 \\
06 & B & 1.461530 & -0.025377 & 1.160801 & 1.259499 & 1.229360 \\
06 & A & 1.487097 & 0.031188 & 1.159354 & 1.293448 & 1.355741 \\
07 & B & 1.483384 & 0.022973 & 1.159627 & 1.265380 & 1.228610 \\
07 & A & 1.493508 & 0.045372 & 1.158365 & 1.280998 & 1.314625 \\
08 & B & 1.463500 & -0.021018 & 1.161653 & 1.342083 & 1.293155 \\
08 & A & 1.493939 & 0.046324 & 1.157274 & 1.310864 & 1.304309 \\
09 & B & 1.474509 & 0.003338 & 1.159803 & 1.280268 & 1.268246 \\
09 & A & 1.491027 & 0.039883 & 1.158652 & 1.291892 & 1.254833 \\
10 & B & 1.479070 & 0.013430 & 1.159930 & 1.335473 & 1.356627 \\
10 & A & 1.480958 & 0.017607 & 1.159871 & 1.351331 & 1.357494 \\
\bottomrule
\end{tabular}
\end{table}

Between-seed price standard deviations are $0.014713$ for B and $0.011766$ for A. Elevation standard deviations are $0.032551$ and $0.026032$. Each run records 100 diagnostic evaluations with 16 episodes each; the requested $10{,}000$-interaction checkpoint spacing is rounded to rollout boundaries, starting at $12{,}288$ interactions. Final outcomes use 64 separate evaluation episodes, and both conditions share seed identifiers and evaluation-stream construction. A mean near the Nash reference does not rule out asymmetry, variation across time, or isolated profitable price cuts. Negative elevation is retained as a valid measured outcome rather than truncated to zero.

Effective feature rank is defined as the exponential Shannon entropy of the normalized eigenvalues of the centered feature covariance matrix. These features are the post-projection, time-augmented representations after $\tanh$ and layer normalization rather than raw GRU hidden states. The reported effective-rank columns are descriptive diagnostics: they do not measure monitoring quality, whether the two firms learned identical representations, or whether a retaliatory punishment rule was encoded. Although the action-independent ablation raises effective rank (from $1.274335$ to $1.338647$ for firm 0 and from $1.269538$ to $1.353578$ for firm 1), it simultaneously changes the marginal input distribution of the public reports. Neither harmful representation collapse under informative monitoring nor a causal effect of feature rank on pricing follows from this comparison.

All 20 rows are included in the reported comparison. A final 64-episode estimate differs slightly from the last 16-episode learning-curve estimate because their evaluation streams differ. The two quantities should not be silently substituted for one another.

%% ============================================================
